\section{Two auxiliary markings and exclusion of bad disks}\label{inc:section}

We impose two conditions on the relative positions of the center and target
endomorphism planes.  The first will be used with Frobenius in
\Cref{nid:frobenius}.  The second rules out proximity to the center at
every place where it is not potentially good.  Both conditions concern
reductions of saturated planes; they impose no primitivity condition on
the particular endomorphisms later chosen to span those planes.

\subsection{Finite-field positions of two planes}

We first give the linear algebra uniformly in the target plane.  Write
$V$ for the split quadratic five-space obtained from a symplectic
four-space by the representation in \Cref{geo:five-space}.  We use the
associated symmetric bilinear form, and call the rank of its restriction
to a subspace the rank of that subspace.
The target is reduced from a saturated rank-two lattice, so it remains
a plane, but its restricted form may have rank zero, one, or two.
The lemma therefore allows an arbitrary target plane $H$.

\begin{lemma}\label{inc:finite-fields}
There is an absolute constant $q_0$ with the following property.  Let
$k$ be a finite field of odd cardinality $q>q_0$, let $H_0\subset V_k$
be a nondegenerate plane, and let $H\subset V_k$ be any plane.  There
are $g_-,g_+\in\Sp_4(k)$ such that
\[
 \dim(H_0+g_-H)=\dim(H_0+g_+H)=4,
\]
the first sum is degenerate, and the second is nondegenerate.
Here the action is the quadratic representation, so the assertion
does not require surjectivity of $\Sp_4(k)\longrightarrow\SO(V_k)$.
\end{lemma}

\begin{proof}
We will use the following elementary observation about rational charts.
The big Bruhat cell of the split group $\Sp_4$ has coordinates
$\mathbb A^4\times\Gm^2\times\mathbb A^4$: multiply the four negative
root subgroups, a diagonal torus element, and the four positive root
subgroups, in a fixed order.  Their matrices, their inverses, and their
actions on $V$ have Laurent-polynomial entries of bounded degrees in
these ten coordinates, independently of $k$.  The same assertion holds
for a big cell of a split parabolic, using its Levi root subgroups and
unipotent radical.  Conjugating a parabolic or inserting a fixed group
element changes coefficients and does not change these degree bounds.

Consequently, an open condition expressed by minors of a bounded-size
matrix in this representation can be tested on either chart by a
polynomial of degree at most an absolute constant $D$.  If the condition
holds at some geometric point, it holds on the corresponding dense
chart: both the group and the parabolic are geometrically irreducible.
Choose a minor that is not identically zero, clear its torus
denominators, and multiply by the torus coordinates.  The resulting
nonzero polynomial has a degree bounded by another absolute constant.
A nonzero polynomial of total degree less than $q$ cannot vanish on
all of $k^n$.  This follows by induction on $n$, applying the univariate
root bound to its nonzero coefficients.  Thus every such geometrically
nonempty open condition has a $k$-point once $q>q_0$, for one constant
$q_0$ that works for all planes under consideration.  No coefficient
height or isometry type enters this bound.

For the nondegenerate sum, work first over $\bar k$.  Choose a
nondegenerate four-space $M$ containing $H_0$, and write
$M=H_0\perp K$ with $\dim K=2$.  Choose bases in which both forms
are $I_2$.  The plane that is the graph of
$A=\diag(a,b):K\longrightarrow H_0$ is disjoint from $H_0$ and has
Gram matrix
\[
 I_2+A^{\mathsf t}A=\diag(1+a^2,1+b^2).
\]
Choices of $a,b\in\bar k$ give each of the ranks $0,1,2$.  Over
$\bar k$, bilinear forms on a two-space are classified by rank, so
there is a graph isometric to $H$, including when $H$ is degenerate.
Witt's extension theorem for subspaces of a nondegenerate space extends
this isometry to an orthogonal transformation of $V_{\bar k}$.
Its determinant can be changed by a reflection in a nonisotropic vector
perpendicular to the graph.  Such a vector exists: that perpendicular
has dimension three, whereas the largest totally isotropic subspace of
$V_{\bar k}$ has dimension two.  We obtain an element of $\SO(V_{\bar
k})$, which lifts to $\Sp_4(\bar k)$ by the spin covering.
Nonvanishing of the Gram determinant of the four concatenated basis
vectors is an open condition and implies their independence.  The
chart argument just proved supplies $g_+\in\Sp_4(k)$ uniformly.

For the degenerate sum, choose an isotropic line $L\subset H_0^\perp$
over $k$.  The perpendicular is a nondegenerate ternary space.  Every
such space over a finite field of odd cardinality is isotropic: after
diagonalization, the two sets
$\{ax^2:x\in k\}$ and $\{-c-by^2:y\in k\}$ each have
$(q+1)/2$ elements and therefore meet, giving an isotropic vector with
third coordinate $1$.  The three-space $H^\perp$ also has an isotropic
line over $k$.  If it is nondegenerate the same argument applies; if
it is degenerate its nonzero radical supplies one.

The group $\Sp_4(k)$ is transitive on the isotropic lines of $V_k$.
Indeed, under the primitive exterior-square realization these lines are
the Pl\"ucker lines of the Lagrangian two-spaces in the symplectic
four-space; symplectic bases show transitivity over $k$ itself.
Move the chosen line in $H^\perp$ to $L$.  After this preliminary
move, $H\subset L^\perp$.  Let $P$ be the stabilizer of $L$ in
$\Sp_4$ and put
\[
 Q=L^\perp/L.
\]
This is a nondegenerate ternary space.  The projection of $H_0$ is
a plane $\overline H_0\subset Q$ because $H_0\cap L=0$.
Over $\bar k$, the Levi action of $P$ on $Q$ includes $\SO(Q)$.
If $H$ contains $L$, its projection is a line.  This line can be
moved outside $\overline H_0$, since the standard representation of
$\SO(Q)$ is irreducible.  Its inverse image then meets $H_0$ trivially.

If $H$ does not contain $L$, its projection is a plane
$\overline H\subset Q$.  Move it so that
$\overline H\ne\overline H_0$.  Such a move exists since a plane
cannot be invariant under the irreducible group $\SO(Q)$.  The two
projected planes now meet in a line $J$.  The unipotent radical permits
an arbitrary change of the lift $Q\longrightarrow L^\perp$ by a
linear map $Q\longrightarrow L$.  Explicitly, choose a hyperbolic
pair $e,f$, with $L=\bar k e$ and $\langle e,f\rangle=1$, and
identify $Q$ with $(\bar k e+\bar k f)^\perp$.  For $z\in Q$ the
orthogonal unipotent transformation is
\[
 e\longmapsto e,\qquad
 n\longmapsto n-\langle n,z\rangle e,\qquad
 f\longmapsto f+z-\tfrac12\langle z,z\rangle e.
\]
These transformations lift to the unipotent radical in $\Sp_4$;
equivalently, in its symplectic realization that radical consists of
the matrices $\left(\begin{smallmatrix}I&Z\\0&I\end{smallmatrix}\right)$
with $Z$ symmetric.  Since the form on $Q$ is nondegenerate,
$n\mapsto-\langle n,z\rangle$ runs through all linear forms on
$Q$.  Choose it so that the lifts of the nonzero vectors of $J$ in
the moved $H$ and in the fixed $H_0$ are distinct.  Then
$H\cap H_0=0$.

We have exhibited a geometric point of $P$ where a four-by-four
minor of the concatenated bases of $H_0$ and the moved $H$ is
nonzero.  The parabolic chart argument gives such a point over $k$
for the same uniform lower bound on $q$.  Their sum, contained in
$L^\perp$ and of dimension four, is exactly $L^\perp$.  Its radical
is $L$, so it is degenerate.  This proves the first assertion and
completes all three possible target-rank cases.
\end{proof}

\subsection{Independent markings on one connected cover}

For a non-CM elliptic-square point $u$, denote its rational Hodge
endomorphism plane by $F_u$.  At an odd prime $\ell$, let
$\Lambda_{u,\ell}$ be the lattice of trace-free self-adjoint operators
on $T_\ell A_u$ from \Cref{geo:five-space}, and set
\begin{equation}\label{inc:saturation}
 \Pi_{u,\ell}=(F_u\otimes\Q_\ell)\cap\Lambda_{u,\ell},
 \qquad
 \overline\Pi_{u,\ell}=\Pi_{u,\ell}/\ell\Pi_{u,\ell}
       \ \subset\Lambda_{u,\ell}/\ell\Lambda_{u,\ell}.
\end{equation}
The quotient $\Lambda_{u,\ell}/\Pi_{u,\ell}$ is torsion free,
so the displayed reduction is an embedded two-space.  A full
symplectic $\ell$-marking identifies its ambient five-space with the
fixed $V_{\mathbb F_\ell}$.  This definition is independent of the
integral pair used to span $F_u$ rationally.

\begin{proposition}\label{inc:markings}
After one fixed finite cover of the curve in \Cref{geo:setup}, there
are distinct fixed odd primes $\ell_1,\ell_2$ and a fixed lift of the
center $t_0$ with the following properties.  Every non-CM elliptic-square
target on the previous curve has a lift $t$ for which, in the respective marked
five-spaces,
\begin{align}
 \overline\Pi_{t_0,\ell_1}+\overline\Pi_{t,\ell_1}
   &\text{ is a degenerate four-space},\label{inc:first}\\
 \overline\Pi_{t_0,\ell_2}+\overline\Pi_{t,\ell_2}
   &\text{ is a nondegenerate four-space}.\label{inc:second}
\end{align}
The prime $\ell_1$ divides none of the fixed Hecke multipliers in
\Cref{geo:hecke-list}.  The prime $\ell_2$ is different from every
residue characteristic at which the fixed center is not potentially
good.  Both center-plane reductions are nondegenerate.  The conditions
remain valid upon simultaneous conjugation of all the data, and the
degree of any chosen lift increases by at most a fixed factor.
\end{proposition}

\begin{proof}
Choose a symplectic integral Betti lattice.  Its geometric monodromy
group $\Gamma\subset\Sp_4(\Z)$ is finitely generated, because the
fundamental group of a smooth algebraic curve is finitely generated,
and is Zariski dense by \Cref{geo:monodromy}.  Strong approximation
for a finitely generated Zariski-dense subgroup of a connected,
simply connected, absolutely almost simple rational group gives
\[
 \Gamma\longrightarrow\Sp_4(\mathbb F_\ell)
 \quad\hbox{surjective for every sufficiently large prime }\ell;
\]
these are precisely the hypotheses of
\cite[Introduction, Theorem, p.~515]{MVW1984} for $\Sp_4$.

Fix $\ell_1>q_0$ outside this theorem's finite exceptional set,
the finitely many primes of degeneracy of the fixed center plane,
and the primes dividing the already fixed multipliers in
\Cref{geo:hecke-list}.  Its reduction kernel
$\Gamma_1$ has finite index in $\Gamma$.  Schreier's theorem gives
finite generation of $\Gamma_1$; its Zariski closure is still
$\Sp_4$, since a connected algebraic group cannot be a finite union
of translates of a proper closed subgroup.  Apply strong approximation
to $\Gamma_1$, and choose $\ell_2>q_0$ distinct from $\ell_1$,
outside its exceptional set, outside the center-plane degeneracy set,
and outside the residue characteristics of all non-potentially-good
places of $A_{t_0}$.  The last set is finite because the fixed abelian
variety has good reduction away from finitely many places.  All these
exclusions use fixed data.

It follows that
\begin{equation}\label{inc:product-reduction}
 \Gamma\longrightarrow
 \Sp_4(\mathbb F_{\ell_1})\times\Sp_4(\mathbb F_{\ell_2})
 \quad\hbox{is surjective}.
\end{equation}
To check this explicitly, first realize a prescribed element in the
first factor by an element of $\Gamma$; correct its second factor
using an element of $\Gamma_1$.  The cover parametrizing the two
full markings with fixed Weil pairings is therefore connected:
its monodromy acts transitively on the product torsor of markings.
After a fixed extension of the field of definition, it is a fixed
geometrically connected finite \'etale cover of our curve, carrying
both markings.  Fix any lift of $t_0$ to it.

For each target, \Cref{inc:finite-fields} prescribes a marking at
$\ell_1$ and, independently, a marking at $\ell_2$ satisfying
\eqref{inc:first} and \eqref{inc:second}.  The full product torsor
lies in the connected cover, so these choices give a point of that
cover.  We continue to write $S,t_0,t$ for the resulting curve,
center, and selected target.  Selection need not be an algebraic rule:
every point in a fiber of a finite morphism of fixed degree has degree
over the original point bounded by that degree.  Subsequent finite
omissions on the curve omit only finitely many points downstairs, as
in \Cref{geo:setup}.

Finally, the planes come from actual rational endomorphisms.  Their
\'etale realizations, intersections with integral Tate lattices,
and induced marked reductions all commute with simultaneous
conjugation.  A conjugation changes the pairing trivialization by the
same similitude on the center and target; on the operator five-space
this is an isometry.  Dimensions and degeneracy are unchanged.
Potential good reduction is also invariant under finite extension,
so adjoining levels and a field of definition creates no new residue
characteristics to be excluded from the choice of $\ell_2$.
\end{proof}

\subsection{A common toric subspace in a strict local model}

We specify the integral models that will also be used in
\Cref{ht:integral-model}.  For each $i=1,2$, forget all levels except
the full $\ell_i$-level and its Weil pairing.  The Siegel datum is
principally polarized, its integral symplectic lattice is self-dual,
and $\ell_i\ge3$ is invertible on
$\Spec\Z[1/\ell_i,\zeta_{\ell_i}]$.  Thus it is the unramified
Siegel case of the integral toroidal construction.  A smooth admissible
projective cone decomposition exists by
\cite[Proposition~7.3.1.4]{Lan2013}; the principal level is neat by
\cite[Remark~1.4.1.9]{Lan2013}.  At this neat level the
construction yields a smooth proper compactification and a
semiabelian extension of the universal abelian scheme; the projective
choice makes the compactification a projective scheme.  We use
\cite[Theorems~6.4.1.1 and~7.3.3.4]{Lan2013}, with the smooth
admissible cone condition of Definition~6.3.3.4, or the Siegel
construction of \cite[Chapter~IV]{FaltingsChai1990}.
In particular we are not identifying an arbitrary toroidal
compactification with a smooth scheme.

After base change to the appropriate ring of integers of our fixed
number field, each base is regular: it is smooth over the regular
Dedekind base with $\ell_i$ inverted.  The semiabelian family is
smooth and separated, of relative dimension two, with geometrically
connected integral fibers.  One of these two models is available at
every residue characteristic, since their levels are coprime.  No
additional common set of primes is inverted here.

We also use the homomorphism extension theorem in its normal-base
form: a homomorphism between two semiabelian schemes on a dense open
of a noetherian normal base extends uniquely to the base
\cite[Chapter~I, Proposition~2.7]{FaltingsChai1990},
\cite[Proposition~3.3.1.5]{Lan2013}.  It applies in particular to
the discrete valuation rings below, after any finite extension needed
to define an endomorphism.  Equivalently, over a trait the semiabelian
extension with abelian generic fiber is the connected part of the
N\'eron model in the semistable case
\cite[\S7.4, Proposition~3 and Corollary~4]{BLR1990}.
We use the extension theorem
directly, so no N\'eron mapping property over a higher-dimensional
ramified base is being assumed.

\begin{lemma}\label{inc:torus}
Let $X$ be one of the regular integral toroidal models just specified,
let $\mathcal G\to X$ be its semiabelian family, and suppose its
interior carries a full $\ell$-marking with $\ell$ invertible on
$X$.  Let $\bar z$ be a geometric point such that
$\mathcal G_{\bar z}$ is a two-dimensional torus.  Any two traits
through $\bar z$ whose generic points lie in the interior have the
same marked two-space
\[
 W_{\bar z}\subset (\Z/\ell\Z)^4
\]
coming from the special torus torsion.  Along each trait it lifts to
a saturated rank-two lattice $\mathcal T\subset T_\ell A$.
Every actual integral endomorphism of the generic abelian variety,
after extension of its field of definition, preserves $\mathcal T$.
If $F\subset\End^0(A)$ is a rational space of endomorphisms, the
reduction of $(F\otimes\Q_\ell)\cap\Lambda_\ell$ therefore
preserves $W_{\bar z}$ as well.
\end{lemma}

\begin{proof}
Take $R=\cO^{\mathrm{sh}}_{X,\bar z}$.  This ring is noetherian
and regular, hence is an integral domain.  The inverse image $U$ of
the moduli interior in $\Spec R$ is a nonempty open in an integral
scheme, hence is connected.  We make this construction on the
regular toroidal base before pulling back to either trait.

For every $n\ge1$, multiplication by $\ell^n$ on the smooth
semiabelian scheme is \'etale: its differential is multiplication by
the invertible integer $\ell^n$.  Its kernel $\mathcal G[\ell^n]$
is consequently separated, quasi-finite and \'etale over $R$.
It need not be finite over $R$.  The finite-part decomposition over
a henselian local base writes it as its finite open-and-closed part
and a part with empty closed fiber.  The finite part is a subgroup;
its formation commutes with local morphisms of henselian local
schemes.  More generally this is the construction of the finite
part and the torus part of quasi-finite flat subgroups in
\cite[\S3.4.1, equations~(3.4.1.1)--(3.4.1.4)]{Lan2013}.

Here the whole special fiber is a torus, so the finite part and the
torus part coincide.  Denote this subgroup by $E_n$.  It is finite
\'etale of rank $\ell^{2n}$, with special fiber
$\mathcal G_{\bar z}[\ell^n]$.  Since $R$ is strictly henselian,
$E_n$ is constant and each of its special points has a unique lift
to an $R$-section.  The multiplication and transition maps among
the $E_n$ are the unique lifts of those on the special torus.  Thus
the system is isomorphic to the standard system
$(\Z/\ell^n\Z)^2$, with its usual reduction transition maps,
after choosing compatible generators in the special torus.

Over $U$, the full $\ell$-marking identifies
$\mathcal G[\ell]_U$ with the constant group $(\Z/\ell\Z)^4$.
Every section of $E_1$ has a constant label under this identification,
because $U$ is connected.  The labels form a subgroup of dimension
two; call it $W_{\bar z}$.  Lift either trait to the strict
henselizations using the chosen geometric specialization.  Compatibility
of the finite part under these local base changes shows that its
generic toric torsion has exactly these labels.  This proves that
the two traits have the same marked subspace.  In particular,
uniqueness is being used over a common strict local base, rather
than separately to compare unrelated trait liftings.

On a trait, passage to the inverse limit of the generic fibers of
$E_n$ gives a rank-two lattice $\mathcal T\subset T_\ell A$.
Its reduction modulo $\ell$ is the embedded two-space $W_{\bar z}$.
It is saturated.  Indeed, lift a basis of $W_{\bar z}$ to compatible
generators of the system $E_n$; the resulting two Tate vectors have
independent reductions in the free rank-four Tate module.  They
extend to a $\Z_\ell$-basis, so their span is a direct summand.
The inverse-limit construction shows that this span is precisely
$\mathcal T$.

An integral endomorphism of the generic fiber extends to the
semiabelian scheme on the trait by the normal-base extension theorem
above.  On the special fiber it preserves the torus and its torsion.
Uniqueness of the lifted torsion sections implies preservation of
each $E_n$ and therefore of $\mathcal T$.  This argument is
compatible with extending the trait, even ramifiedly.  Rational
linear combinations preserve $\mathcal T_{\Q_\ell}$.
Since $\mathcal T$ is saturated,
\[
 \mathcal T=\mathcal T_{\Q_\ell}\cap T_\ell A.
\]
Every integral Tate operator lying in their rational span preserves
this intersection.  Reducing that assertion modulo $\ell$ proves
the last statement, in particular for the full saturation of an
endomorphism plane even if its chosen spanning pair becomes
dependent modulo~$\ell$.
\end{proof}

\Needspace{8\baselineskip}
\subsection{Excluding proximity at a bad center place}

\begin{lemma}\label{inc:stabilizer}
Let $V_4$ be a symplectic four-space over a field of characteristic
different from two.  No nondegenerate four-space of trace-free
symplectic-self-adjoint operators can preserve a two-space
$W\subset V_4$.
\end{lemma}

\begin{proof}
Write $\psi$ for the symplectic form. Its restriction to $W$ has rank
two or zero. In the first
case $V_4=W\perp W^\perp$.  A self-adjoint operator preserving
$W$ also preserves $W^\perp$, since
$\psi(Xv,w)=\psi(v,Xw)$.
A self-adjoint operator on a symplectic two-space is scalar, as is
seen by solving $D^{\mathsf t}J=JD$ for a two-by-two matrix.
Our operator is consequently $aI_W\oplus bI_{W^\perp}$, and its
trace-zero condition is $2a+2b=0$.  The space of such operators has
dimension one.

In the rank-zero case $W$ is Lagrangian.  Choose it as the first
Lagrangian in a symplectic basis, with form matrix
$\left(\begin{smallmatrix}0&I\\-I&0\end{smallmatrix}\right)$.
The self-adjoint trace-free operators preserving $W$ are exactly
\begin{equation}\label{inc:lagrangian-stabilizer}
 \left\{
 X(D,J)=\begin{pmatrix}D&J\\0&D^{\mathsf t}\end{pmatrix}:
 \tr D=0,\quad J^{\mathsf t}=-J
 \right\}.
\end{equation}
This is a four-space, since $\dim\mathfrak{sl}_2=3$ and the
skew-symmetric two-by-two matrices have dimension one.  Its trace
pairing is
\[
 \Tr\bigl(X(D,J)X(D',J')\bigr)=2\tr(DD').
\]
The trace pairing on $\mathfrak{sl}_2$ is nondegenerate in odd
characteristic: in coordinates
$D=\left(\begin{smallmatrix}a&b\\c&-a\end{smallmatrix}\right)$
it pairs $(a,b,c)$ with $(a',b',c')$ as
$2aa'+bc'+cb'$.  Hence the radical of
\eqref{inc:lagrangian-stabilizer} is exactly the one-dimensional
$J$-space.  Any four-space of operators contained in this stabilizer
equals it and is degenerate.  Both cases give the required
contradiction.
\end{proof}

\begin{proposition}\label{inc:bad-disks}
Fix a finite place of the field of the center at which $A_{t_0}$
is not potentially good.  There is a sufficiently small disk about
the actual point $t_0$ in the completed local curve such that no
selected target from \Cref{inc:markings} lies in this disk.  The
radius is independent of the finite, possibly ramified extension
over which a target is defined.  The same conclusion holds for
every conjugate of the fixed data, with a fixed finite collection
of radii.
\end{proposition}

\begin{proof}
We first identify the special group at the center.  After a finite
local extension the elliptic curve $E$ in an isogeny
$A_{t_0}\sim E^2$ has semistable reduction.  One may obtain this
extension by choosing a full elliptic level invertible at the place
and applying the same proper toroidal construction in dimension one.
An elliptic curve with
semistable reduction is either good or multiplicative.  Toric rank
is invariant under isogeny and is additive on products: extend an
isogeny and a quasi-inverse, with composites multiplication by an
integer, to the semiabelian models.  On special fibers their
restrictions map tori into tori, since a homomorphism from a torus
to an abelian variety is zero.  The two composites on the tori
are multiplication by that integer, so these restrictions are
isogenies and the torus dimensions agree.  Thus the potential
toric rank of $A_{t_0}$ is zero or two.  Rank zero in the
semistable model is good reduction; in that case $A_{t_0}$ would
be potentially good.  Our hypothesis therefore gives rank two.

Use the full $\ell_2$-level toroidal model, which is available at
this residue characteristic by \Cref{inc:markings}.  Properness
extends the center point to a section over the valuation ring,
after a fixed local extension if needed.  Its pullback semiabelian
scheme is the semistable model of the center, by uniqueness of
semiabelian extension over a normal trait.  Its geometric special
fiber is therefore a two-dimensional torus.  Denote the
specialization in the toroidal base by $\bar z$.

We show that a disk of fixed radius forces this very same
specialization for every target.  Choose an affine open of the
toroidal scheme containing $\bar z$, and enlarge the fixed local
field so that the center section and the required residue point
are defined.  The generic image of $t_0$ lies in this affine open.
Finitely many affine coordinate generators pull back to rational
functions on the curve that are regular at $t_0$; shrink its
parameter neighborhood also to exclude their denominators and
the complement of this affine open.  Their values at $t_0$ are
integral. Choose a local parameter $x$ at $t_0$ with $x(t_0)=0$.
Continuity gives a real $r>0$ such that
\[
 |x(t)|<r
 \quad\Longrightarrow\quad
 |f_j(t)-f_j(t_0)|<1
 \qquad\hbox{for every coordinate generator } f_j.
\]
The assertion uses the branch of $x$ through $t_0$.  Each $f_j$
has a convergent expansion on a fixed smaller disk, so the same
inequalities hold in the fixed algebraic closure of the local
field, for points in every finite extension.  They do not involve
a valuation normalization by a varying ramification index.
The values $f_j(t)$ are integral and satisfy the equations of
the chosen affine chart, and have the center's residues.
They consequently define an integral section with specialization
$\bar z$.  Separatedness identifies its generic point with
the given moduli point of $t$.

If a selected target lay in this disk, the two traits would
therefore satisfy \Cref{inc:torus} with $\ell=\ell_2$.
Both $\overline\Pi_{t_0,\ell_2}$ and
$\overline\Pi_{t,\ell_2}$ would preserve the same two-space
$W_{\bar z}$ in the marked symplectic four-space.  Their sum
is a nondegenerate four-space by \eqref{inc:second}, in direct
contradiction to \Cref{inc:stabilizer}.  This proves exclusion
of the target point from the disk itself.

The center has only finitely many non-potentially-good places,
and the fixed data have only finitely many conjugates.  Apply
the argument at each of them.  Simultaneous conjugation preserves
\eqref{inc:second}, so it gives finitely many fixed radii with
the asserted properties.  All of this argument uses integral
models and prime-to-characteristic torsion.  No comparison with
de Rham horizontal transport at a bad center is needed.
\end{proof}
